\SourceLeaf{337}%
mangelung aller andern Kenntniss \(T^0\) annehmen mag), die vier Formeln berechnen
\[
\begin{aligned}
\beta&=(2)r\cos T\\
B&=B^0-\tfrac12\beta\\
\tau&=(1)r\sin T\operatorname{tang}B\\
T&=T^0-\tfrac12\tau
\end{aligned}
\]
und zwar wird man den Werth von (2), der aus der Tafel mit dem Argument \(B\) entnommen werden sollte, das erstemal mit dem Argument \(B^0\) entnehmen können, wenn man nicht durch Schätzung einen schon mehr genäherten Werth von \(B\) anticipiren zu können glaubt; den Werth von (1) nimmt man aus der Tafel mit dem eben gefundenen Werthe von \(B\).

Dieselbe Rechnung wiederholt man mit dem durch die vierte Gleichung gefundenen Werthe von \(T\), indem man (1) und (2) mit dem schon verbesserten \(B\) aus der Tafel entlehnt; und so macht man nöthigenfalls eine abermalige Wiederholung, bis das Resultat zum Stehen kommt, d.\,i. bis man durch die vierte Formel denselben Werth von \(T\) wiedererhält, von dem man zuletzt ausgegangen war. Zu allen diesen Rechnungen wird man nur fünfzifrige Logarithmen verwenden.

Bei den weitern Wiederholungen wird man die Rechnung mit siebenzifrigen Logarithmen führen, die logarithmischen Correctionen von \(\log\tau\) und \(\log\beta\) mit zuziehen, und \(B=B^0-\tfrac12b,\ T=T^0-\tfrac12t\) setzen. Erst wenn auch diese Rechnung stehende Resultate gegeben hat, wird man auch \(\lambda\) und \(l\) nach den am Schluss des 33. Art. gegebenen Formeln berechnen. Zur Erläuterung dieser Vorschriften mögen hier die Hauptmomente eines Beispiels stehen, welches eben so wie oben Art. 20 bei der sphärischen Rechnung von der Dreiecksseite Brocken-Inselsberg hergenommen ist.

Bei der ellipsoidischen Rechnung ist die Breite des Brockens
\[
=51^0\ 48'\ 1''\ 9294=B^0
\]
das Azimuth der Seite Brocken-Inselsberg
\[
=5^0\ 42'\ 21''\ 7699=T^0
\]
Der Logarithm der Dreiecksseite in Toisen ist bis auf die siebente Decimale derselbe wie in der conformen Darstellung auf der Kugelfläche, nemlich \(=4{,}7353929\), folglich in der unsrer Hülfstafel zum Grunde liegenden Einheit \(\log r=5{,}0251757\).
